\documentclass{article}

\usepackage{amsmath}
\usepackage{booktabs}
\usepackage{enumitem}
\usepackage[most]{tcolorbox}

\newtcolorbox{eqbox}{colback=blue!6!white, colframe=blue!45!black,
  boxrule=0.8pt, sharp corners, left=2pt, right=2pt, top=2pt, bottom=2pt}
\newtcolorbox{whybox}{colback=black!4!white, colframe=black!4!white,
  borderline west={2pt}{0pt}{black!30}, sharp corners,
  left=6pt, right=6pt, top=4pt, bottom=4pt}

% Compact variable-definition list placed under each equation
\newlist{vardefs}{itemize}{1}
\setlist[vardefs]{label=--, nosep, leftmargin=1.6em, topsep=2pt}

\begin{document}

\section{Strategy and Experimental Design}

\subsection{Hedging Strategy: Static No-Trade Band}

The strategy monitors the book's net delta $\Delta(t)$ on a 5-minute grid and keeps a
static no-trade band around the target $\Delta_{\text{target}} = 0$:

\begin{eqbox}
\begin{equation}
-H \le \Delta(t) \le H.
\end{equation}
\end{eqbox}

\begin{vardefs}
\item $\Delta(t)$: net delta of the option book at time $t$, measured in ES futures
  contract equivalents (positive = long the market);
\item $H$: half-width of the no-trade band, in ES contract equivalents.
\end{vardefs}

While $\Delta(t)$ remains inside the band, the system does nothing and no execution cost
is incurred. A hedge is triggered only when the breach clears the band by a buffer
$\epsilon$:

\begin{eqbox}
\begin{equation}
\lvert \Delta(t) \rvert > H(1+\epsilon),
\label{eq:trigger}
\end{equation}
\end{eqbox}

\begin{vardefs}
\item $\epsilon$: breach buffer, expressed as a fraction of $H$ (e.g.\ $\epsilon = 5\%$
  requires the delta to clear the band by an extra 5\% before hedging).
\end{vardefs}

Once triggered, the parent hedge order returns the book to flat:

\begin{eqbox}
\begin{equation}
Q_t = -\Delta(t).
\end{equation}
\end{eqbox}

\begin{vardefs}
\item $Q_t$: signed parent hedge order in ES futures contracts (a long book sells
  futures, a short book buys futures).
\end{vardefs}

\textbf{Post-hedge cooldown.} After a hedge schedule completes, the trigger in
Equation~\eqref{eq:trigger} is suppressed for a minimum interval $T_{\text{cooldown}}$,
unless the breach is large ($\lvert \Delta(t) \rvert > 2H$). This prevents the system from
immediately re-triggering a fresh schedule on the next small delta arrival. At market
close, any remaining delta is forcibly flattened, so every strategy configuration
satisfies the delta-flat overnight constraint $\Delta(T) = 0$ by construction.

\subsection{Execution: TWAP and VWAP Schedules}

Once triggered, the parent order $Q_t$ is split into child orders over an execution
horizon $\tau$, i.e.\ over $N = \tau / 5\,\text{min}$ consecutive bars. Under TWAP the
child orders are equal-sized; under VWAP they are proportional to the ES futures volume
observed in each bar of the horizon:

\begin{eqbox}
\begin{equation}
q_k^{\text{TWAP}} = \frac{Q_t}{N},
\qquad
q_k^{\text{VWAP}} = Q_t\,\hat{w}_k,
\qquad
\sum_{k=1}^{N} \hat{w}_k = 1.
\end{equation}
\end{eqbox}

\begin{vardefs}
\item $q_k$: signed child order executed in bar $k$ of the schedule, in ES contracts;
\item $\tau$: execution horizon after the trigger (minutes);
\item $N$: number of 5-minute bars in the horizon, $N = \tau / 5\,\text{min}$;
\item $\hat{w}_k$: bar $k$'s share of ES futures volume over the horizon, with
  $\sum_k \hat{w}_k = 1$.
\end{vardefs}

If a new trigger fires while a schedule is still running, the old schedule is cancelled
and replaced by a fresh one sized to the current residual delta.

\subsection{Units and Book-Delta Measurement}

All delta and hedge quantities are measured in ES futures contract equivalents. For each
executed ES option trade $i$, its book-delta contribution is

\begin{eqbox}
\begin{equation}
\Delta_{\text{new},i} = s_i \cdot n_i \cdot \delta_i.
\label{eq:arrival}
\end{equation}
\end{eqbox}

\begin{vardefs}
\item $\Delta_{\text{new},i}$: book-delta contribution of option trade $i$, in ES
  contract equivalents;
\item $s_i \in \{+1,-1\}$: trade direction sign ($+1$ = client buy, $-1$ = client sell);
\item $n_i$: number of option contracts in trade $i$;
\item $\delta_i$: per-contract delta of the traded option, $\delta_i \in [-1, 1]$.
\end{vardefs}

The ES futures multiplier $M_h = 50$ is \emph{not} applied in
Equation~\eqref{eq:arrival}: an ES option's delta is already expressed in ES futures
contract equivalents, and $M_h$ enters only when converting futures price moves or
execution costs into dollars.

Because the hedge instrument is an ES \emph{future}, option deltas are computed with the
Black-76 futures-option delta rather than the equity Black--Scholes delta. At bar $t_j$,

\begin{eqbox}
\begin{equation}
\delta_i(t_j) =
\begin{cases}
N(d_{1,i}(t_j)), & \text{call},\\[2pt]
N(d_{1,i}(t_j)) - 1, & \text{put},
\end{cases}
\qquad
d_{1,i}(t_j) = \frac{\ln(F_j/K_i) + \tfrac{1}{2}\sigma^2 \tau_i(t_j)}{\sigma\sqrt{\tau_i(t_j)}}.
\label{eq:delta}
\end{equation}
\end{eqbox}

\begin{vardefs}
\item $\delta_i(t_j)$: delta of option $i$ evaluated at bar $t_j$;
\item $N(\cdot)$: standard normal cumulative distribution function;
\item $F_j$: ES futures price at bar $t_j$;
\item $K_i$: strike price of option $i$;
\item $\tau_i(t_j)$: time to expiry of option $i$ at bar $t_j$, in years;
\item $\sigma$: flat volatility input, the trading day's 1M BVOL.
\end{vardefs}

The $e^{-r\tau}$ discount factor of the fully discounted Black-76 delta is omitted; at
intraday hedging horizons its effect on delta is negligible.

\textbf{Continuous revaluation.} The book delta at bar $t_j$ is recomputed from scratch
by repricing \emph{every} live option position at the current futures price,

\begin{eqbox}
\begin{equation}
\Delta_{\text{book}}(t_j) = \sum_{i:\,t_i \le t_j} s_i \cdot n_i \cdot \delta_i(t_j),
\qquad
\Delta(t_j) = \Delta_{\text{book}}(t_j) + \Delta_{\text{hedge}}(t_j).
\end{equation}
\end{eqbox}

\begin{vardefs}
\item $\Delta_{\text{book}}(t_j)$: delta of all option positions accumulated up to bar
  $t_j$, repriced at $F_j$;
\item $t_i$: execution timestamp of option trade $i$ (the sum runs over all trades
  executed at or before $t_j$);
\item $\Delta_{\text{hedge}}(t_j)$: cumulative sum of executed hedge child orders up to
  bar $t_j$;
\item $\Delta(t_j)$: residual book delta after hedging --- the quantity the band in
  Equation~\eqref{eq:trigger} monitors.
\end{vardefs}

\begin{whybox}
\textbf{Why revaluation matters.} A cumulative sum of deltas frozen at each trade's own
execution price only moves when new trades arrive. Repricing the live book each bar
captures the gamma-driven drift of positions already on the book, which is a material
share of intraday delta movement and is required to reproduce the continuously drifting
zero-hedge delta path.
\end{whybox}

\subsection{Execution-Cost and Delta-Risk Measurement}

Actual hedge fills are not available, so execution cost is estimated with a
market-data-based proxy combining spread cost and an ADV-adjusted impact penalty:

\begin{eqbox}
\begin{equation}
C_{\text{exec}} = \sum_k M_h \lvert q_k \rvert
\left[ \frac{s}{2} + \alpha S_k \left( \frac{\lvert q_k \rvert}{ADV_k} \right)^{\!\beta} \right].
\label{eq:cost}
\end{equation}
\end{eqbox}

\begin{vardefs}
\item $C_{\text{exec}}$: total execution cost in dollars, summed over all hedge child
  orders $k$ (the forced end-of-day flattening trade included);
\item $M_h = 50$: ES futures multiplier (dollars per index point per contract);
\item $q_k$: signed child order size in ES contracts;
\item $s/2$: half bid--ask spread in index points, proxied by half of one ES tick
  ($0.125$), since intraday quotes are not in the 5-minute bar data and front-month ES
  trades at the minimum tick spread almost continuously;
\item $\alpha$: market-impact level parameter (base case $0.05$);
\item $S_k$: ES futures price at the bar where child order $k$ executes;
\item $ADV_k$: 20-day average daily ES futures volume, in contracts;
\item $\beta$: market-impact curvature parameter (base case $0.5$); the ratio
  $\lvert q_k \rvert / ADV_k$ is the dimensionless participation rate.
\end{vardefs}

Delta risk is measured from the realized residual-delta P\&L of the simulated hedge
path:

\begin{eqbox}
\begin{equation}
\ell^{\Delta}_{j+1} = M_h\,\Delta(t_j)\,(S_{j+1} - S_j),
\qquad
C_{\text{risk}} = \sqrt{ \frac{1}{M} \sum_{j=0}^{M-1} \left( \ell^{\Delta}_{j+1} \right)^{2} }.
\end{equation}
\end{eqbox}

\begin{vardefs}
\item $\ell^{\Delta}_{j+1}$: realized P\&L (dollars) from carrying residual delta
  $\Delta(t_j)$ through the price move from bar $j$ to bar $j+1$;
\item $S_j$: ES futures price at bar $t_j$;
\item $M$: total number of 5-minute bars in the backtest;
\item $C_{\text{risk}}$: delta risk, the root mean square of realized residual-delta
  P\&L across all bars.
\end{vardefs}

\subsection{Benchmarks and Standardized Objective}

Two benchmark policies bound the attainable cost--risk region
(Table~\ref{tab:benchmarks}). \textbf{Zero Hedge} never rebalances intraday and flattens
only at the close: it is the lower bound on execution cost and the upper bound on delta
risk. \textbf{Full Hedge} re-flattens the book at every 5-minute bar --- offsetting both
new arrivals and the gamma drift of existing positions --- and is the upper bound on cost
with delta risk identically zero.

\begin{table}[htbp]
\centering
\small
\begin{tabular}{lll}
\toprule
\textbf{Benchmark} & \textbf{Description} & \textbf{Purpose} \\
\midrule
Zero Hedge & Flatten at close only & Cost floor, risk ceiling \\
Full Hedge & Re-flatten every 5-minute bar & Risk floor, cost ceiling \\
\bottomrule
\end{tabular}
\caption{Benchmark policies used for normalization.}
\label{tab:benchmarks}
\end{table}

Each strategy configuration is scored by normalizing its cost against Full Hedge and its
risk against Zero Hedge, so both terms are dimensionless and $O(1)$:

\begin{eqbox}
\begin{equation}
\widetilde{C}_{\text{exec}}(\theta) = \frac{C_{\text{exec}}(\theta)}{C_{\text{exec}}(\text{Full Hedge})},
\qquad
\widetilde{C}_{\text{risk}}(\theta) = \frac{C_{\text{risk}}(\theta)}{C_{\text{risk}}(\text{Zero Hedge})},
\end{equation}
\end{eqbox}

\begin{eqbox}
\begin{equation}
L_{\text{std}}(\theta) = \widetilde{C}_{\text{exec}}(\theta)
+ \lambda_{\text{risk}}\,\widetilde{C}_{\text{risk}}(\theta),
\qquad
\theta^{*} = \arg\min_{\theta \in \Theta} L_{\text{std}}(\theta).
\end{equation}
\end{eqbox}

\begin{vardefs}
\item $\theta = (H, \epsilon, T_{\text{cooldown}}, e, \tau)$: one strategy configuration
  --- band width, breach buffer, cooldown, execution method ($e \in \{\text{TWAP},
  \text{VWAP}\}$), and execution horizon;
\item $\widetilde{C}_{\text{exec}}(\theta)$: execution cost of $\theta$ as a fraction of
  Full Hedge's cost;
\item $\widetilde{C}_{\text{risk}}(\theta)$: delta risk of $\theta$ as a fraction of
  Zero Hedge's risk;
\item $\lambda_{\text{risk}}$: weight on normalized delta risk (base case $1$, i.e.\
  equal weight on cost and risk; swept as a sensitivity parameter);
\item $L_{\text{std}}(\theta)$: standardized loss to be minimized;
\item $\Theta$: the full parameter grid; $\theta^{*}$ is the selected configuration.
\end{vardefs}

\subsection{Experimental Design}

\textbf{Data.} The backtest covers 102 trading days (January 2 -- May 29, 2026). Delta
arrivals come from a simulated tape of 32{,}127 executed ES option trades calibrated to
Databento ES option prints; market inputs are 8{,}364 five-minute ES futures bars with
volume, 20-day ADV, 1M BVOL, and 1M SOFR OIS.

\textbf{Empirical calibration of the band grid.} The band width $H$ is a parameter to be
selected by backtesting, and its scale must match the book it polices. Under the unit
convention of Equation~\eqref{eq:arrival}, the realized book delta has a median daily
maximum of $\lvert \Delta \rvert \approx 85$ contracts (95th percentile $\approx 257$,
maximum $609$). The band grid is therefore chosen to span this empirical range,

\begin{eqbox}
\begin{equation}
H \in \{25,\ 50,\ 100,\ 200\} \ \text{ES contract equivalents},
\end{equation}
\end{eqbox}

so that the narrowest band binds on most days and the widest binds only on tail days. As
limiting cases, $H \to 0$ recovers Full Hedge and $H \to \infty$ recovers Zero Hedge.

\begin{whybox}
\textbf{Why not the illustrative grid.} Band values quoted in earlier drafts (e.g.\
$H \in \{2{,}500,\dots,20{,}000\}$) sit one to two orders of magnitude above this book's
realized delta scale; at that scale the band never triggers and every configuration
collapses to Zero Hedge. Calibrating $H$ to the empirical delta distribution is what
makes the grid search informative.
\end{whybox}

\textbf{Parameter grid.} The full strategy grid crosses every dimension of $\theta$
(Table~\ref{tab:grid}), giving $4 \times 3 \times 3 \times 3 \times 2 = 216$
configurations. Each is simulated over all 102 days with the same precomputed book-delta
paths, so differences across configurations reflect the hedging rule alone.

\begin{table}[htbp]
\centering
\small
\setlength{\tabcolsep}{5pt}
\begin{tabular}{llll}
\toprule
\textbf{Parameter} & \textbf{Role} & \textbf{Grid} & \textbf{Selection} \\
\midrule
$H$ & No-trade band width & $\{25, 50, 100, 200\}$ contracts & Grid search \\
$\epsilon$ & Breach buffer & $\{0\%, 5\%, 10\%\}$ & Grid search \\
$T_{\text{cooldown}}$ & Post-hedge cooldown & $\{0, 15, 30\}$ min & Grid search \\
$\tau$ & Execution horizon & $\{15, 30, 60\}$ min & Grid search \\
$e$ & Execution method & $\{\text{TWAP}, \text{VWAP}\}$ & Comparison \\
\midrule
$\alpha$ & Impact level & $0.05 \times \{0.5, 1, 2, 4\}$ & Sweep \\
$s/2$ & Half-spread proxy & $0.125 \times \{0.5, 1, 2, 4\}$ pts & Sweep \\
$\lambda_{\text{risk}}$ & Risk weight & $\{0.1, 0.25, 0.5, 1, 2, 5, 10\}$ & Sweep \\
\bottomrule
\end{tabular}
\caption{Strategy parameters (top) and model-assumption parameters (bottom). $\beta$ is
held at $0.5$: rescaling the impact exponent requires per-order participation rates that
are not retained at the aggregate grain of the results.}
\label{tab:grid}
\end{table}

\textbf{Evaluation.} Each configuration is placed in normalized (cost, risk) space
against the two benchmarks, and the Pareto frontier of non-dominated configurations is
traced. The selected strategy is $\theta^{*}$, the minimizer of
$L_{\text{std}}$; its robustness is then checked under the $\alpha$, spread, and
$\lambda_{\text{risk}}$ sweeps, and its behavior is illustrated with an intraday
walk-through (book-delta path and cumulative delta P\&L) on the busiest simulated
trading day.

\end{document}
